\expandafter\def\csname HistoricalNoteText9\endcsname{When $S>\BareRoot{\frac12}$, the first condition implies the second; but when $S<\BareRoot{\frac12}$, the second implies the first.}

\expandafter\def\csname HistoricalNoteText10\endcsname{It certainly seems to have been demonstrated well enough that an algebraic curve can neither suddenly break off anywhere (as happens, for example, with the transcendental curve whose equation is $y=\frac1{\mathrm{log}\,x}$), nor, after infinitely many windings, as it were lose itself at some point (like the logarithmic spiral); and, as far as I know, no one has raised a doubt on this matter. Nevertheless, if anyone demands it, I shall undertake to give a proof free from all doubt on another occasion. In the present case, however, it is clear that if some branch, for example 2, never left the circle anywhere (Fig. 3), you could enter the circle between 0 and 2, then go around the whole of this branch (which would have to lose itself in the region of the circle), and finally leave the circle again between 2 and 4, without encountering the first line anywhere along the entire path. That this is absurd is clear from the fact that, at the point where you entered the circle, the first surface was above you, and at your exit it was below you. You must therefore necessarily have encountered the first surface itself somewhere, or a point of the first line. --- This reasoning, which rests on principles of the geometry of position that are no less valid than those of the geometry of magnitude, implies only that, if you enter the circle along some branch of the first line, you can leave it again at another place while always remaining on the first line; it does not imply that your path is a continuous line in the sense accepted in higher geometry. But here it is enough for the path to be a continuous line in the ordinary sense, that is, nowhere interrupted but connected everywhere.}
